\section*{Discussion}

\textbf{Key Findings and Implications.} Our results demonstrate that physician-specific adaptation is both necessary and effective: at the 3B scale, the gap between the base model (BLEU-4: 5.1, Personalization: 2.1) and individualized LoRA adaptation (BLEU-4: 10.4, Personalization: 4.2) confirms that generic models cannot capture individual practice styles. SCF further enables practical deployment by preserving individual adapter performance post-merging, allowing hundreds of physicians to be served through a single inference endpoint---a critical advance for resource-constrained hospital settings. The DAR router supports both identity-based and query-based deployment scenarios with minimal latency overhead.


\textbf{Limitations.} Several limitations warrant discussion~\cite{crossBiasMedicalAI2024Nian11Yue7Ri}. First, IHDPD originates from a single internet hospital; cross-institutional validation across diverse healthcare systems is needed. Second, our approach captures stylistic preferences rather than clinical reasoning, limiting applicability where diagnostic logic is the key differentiator. Third, under suboptimal SCF trim fractions, a hybrid deployment where outlier physicians retain individual adapters may bridge the remaining gap. Fourth, LLM-as-a-Judge assessments should be complemented by independent physician evaluators in future work.


\textbf{Ethical Considerations.} Ethically, personalized LLM deployment must preserve physician autonomy~\cite{zhangEthicsGovernanceTrustworthy2023}. The final decision to adopt or reject AI-generated suggestions remains with the clinician, and institutional policies should mandate appropriate-use training to prevent automation bias. We note that physician-specific personalization may actually reduce uncritical adoption compared to generic outputs, as physicians can more readily evaluate suggestions framed in their own clinical language. 
% Clear governance frameworks for accountability remain essential.

\textbf{Future Work.} Promising extensions include multi-modal personalization (e.g., diagnostic imaging reports), continuous LoRA updating as physicians accumulate new consultations, and DAR enhancement via reinforcement learning from physician adoption feedback.



